\documentclass[10pt]{article}
\usepackage[margin=2.2cm]{geometry}
\usepackage{amsmath}
\usepackage{amssymb}
\usepackage{booktabs}
\usepackage{caption}
\captionsetup{font=small}
\begin{document}

\begin{center}
{\Large\bfseries Supplementary Material}\\[4pt]
{\large\bfseries Adapter Scaling Trade-off and Timestep-Conditioned Scheduling in Multi-view Diffusion Texture Generation}\\[4pt]
{\large FAC controlled re-examination: full dose--response study}
\end{center}

\section{S1.\ \ Protocol}
All variants are trained and evaluated under the strictly paired protocol of Sec.~4.6 of the main manuscript: the base model, VAE, and GeoTex-Adapter remain frozen; only the LTAG/GSG/FSC controllers (roughly $6\times10^{3}$ parameters) are trained with the same enhanced denoising objective as the base adapter (a latent noise-prediction MSE with $1.7\times$ foreground and $1.3\times$ edge spatial reweighting plus a latent-SSIM term; AdamW, peak learning rate $10^{-4}$ with 100 warm-up steps), on the adapter's training pool, which is strictly disjoint from all evaluation objects; evaluation uses 300 objects, the Euler discrete scheduler with 50 steps, seed 42, and identical evaluation code for every variant. Deficits are paired per object against the training-free TCAS baseline ($C3=(1.25,2.50,1.25)$).

\section{S2.\ \ Dose--response over seven trained configurations}

\begin{table}[h]
\centering\small
\caption{FAC controlled re-examination: seven trained configurations, all evaluated under the strictly paired protocol. FG-PSNR and FG-SSIM are foreground-masked absolute values; $\Delta$FG-PSNR is the paired mean difference relative to the training-free TCAS baseline; the win rate is the fraction of objects on which the variant improves foreground PSNR. The untrained warm-start control (12-object paired subset, gate fixed at the exact C3 envelope) is reported in the text: $\Delta$FG-PSNR $=-0.18$~dB, 95\% CI $[-0.42,+0.05]$, statistically indistinguishable from the baseline.}
\label{tab:fac_dose}
\begin{tabular}{llcccc}
\toprule
Configuration & Temporal gate & FG-PSNR & FG-SSIM & $\Delta$FG-PSNR & Win rate \\
\midrule
TCAS baseline (C3) & fixed C3 & 9.88 & 0.3681 & -- & -- \\
\midrule
LTAG only & joint, unconstrained & 7.79 & 0.3487 & $-2.08$ & 4/300 \\
LTAG+GSG & joint, unconstrained & 7.03 & 0.3163 & $-2.84$ & 6/300 \\
Full FAC & joint, unconstrained & 6.94 & 0.3204 & $-2.94$ & 5/300 \\
Two-stage + GSG/FSC & frozen at C3 & 7.32 & 0.3074 & $-2.56$ & 4/300 \\
\ \ + envelope penalty $\lambda{=}1$ & joint, penalized to C3 & 7.14 & 0.3021 & $-2.74$ & 2/300 \\
\ \ + fidelity weighting & frozen at C3 & 7.51 & 0.3110 & $-2.37$ & 5/300 \\
\ \ + gentle schedule & frozen at C3 & 8.77 & 0.4175 & $-1.10$ & 48/300 \\
\bottomrule
\end{tabular}
\end{table}

\section{S3.\ \ Reading of the dose--response pattern}

Three observations follow from Table~\ref{tab:fac_dose}. First, the warm-start control shows that when the temporal gate is held at the exact C3 envelope and not trained, the extension is statistically indistinguishable from the training-free schedule (paired $\Delta=-0.18$~dB, 95\% CI $[-0.42,+0.05]$, crossing zero); the evaluation path is therefore validated. Second, the deficit grows monotonically with the amount of controller training under the denoising objective: zero steps preserve the baseline, gentle training (500 steps at a fourfold-reduced learning rate) costs $1.10$~dB, and full training costs up to $2.94$~dB. Third, envelope preservation is necessary but not sufficient: even with the gate frozen at C3, training the spatial and frequency controllers degrades fidelity, and the gentle variant merely trades signal fidelity for a slightly higher structural statistic (FG-SSIM $+0.049$ at FG-PSNR $-1.10$~dB) rather than achieving a genuine improvement.

The mechanism is consistent with the controllers optimizing the denoising objective at the expense of foreground fidelity: gradient training on the same loss as the base adapter has no term that protects the fixed schedule's balance, so the learned modulation migrates away from the operating point that makes TCAS effective. We therefore identify, as future work, training signals that do not trade away foreground fidelity (fidelity-weighted or structure-aware objectives), envelope-preserving parameterizations and constraints, and post-hoc rather than gradient-learned gates. With the tested lightweight controllers, the training-free TCAS remains the main and sole method of the paper, and all its conclusions are unaffected by this negative result.

\section{S4.\ \ Reproducibility details}

Table~\ref{tab:repro} consolidates the data, checkpoint, and inference details requested by the reviewers. All identifiers refer to the internal object list of the 300-object evaluation pool; checkpoint files and evaluation code are made available upon reasonable request.

\begin{table}[h]
\centering\small
\caption{Consolidated reproducibility table.}
\label{tab:repro}
\begin{tabular}{p{4.3cm}p{10.8cm}}
\toprule
Item & Setting \\
\midrule
Data source & Objaverse 3D asset collection (including ABO subset assets), rendered offline \\
Rendering & Blender (Cycles engine), 17 views, $512\times512$, white background \\
Training pool (main adapter) & 1{,}118 rendered objects (all main-text tables) \\
Training pool (strong-residual instance) & 1{,}706 rendered objects from the same source; used by the strong-residual, per-layer-capped instance of Sec.~4.7, the FAC extension, and this supplementary study \\
Evaluation pool & 300 rendered objects in a fixed ordered list; probe $=$ first 24 identifiers (obj$_{0000}$--obj$_{0023}$), held-out test $=$ remaining 276 (obj$_{0024}$--obj$_{0299}$) \\
Train--evaluation overlap & Zero (training and evaluation identifier sets verified to intersect in zero objects) \\
Base pipeline & MVPainter-style geometry-controlled multi-view diffusion model; joint six-view latent UNet (pipeline class \texttt{MVPainter\_Pipeline}, diffusers 0.20.0) \\
Base checkpoint & Local MVPainter snapshot (fine-tuned UNet, VAE, and dual CLIP vision encoders; loaded frozen) \\
Adapter checkpoint (main tables) & GeoTex-Adapter \texttt{refattn\_v1} (\texttt{geotex\_step\_0002000.pt}); requested scale applied multiplicatively at every injected layer \\
Adapter checkpoint (Sec.~4.7 / FAC / supp.) & \texttt{geotex\_v2\_ema\_final.pt} (MD5 \texttt{a74cc1d16cb473a95022bdcbc58e1b22}, 8.86~MB), applied with per-layer caps deep 3.0 / middle 3.5 / shallow 0.8 \\
Resolution & $512\times512$ for rendering, generation, and evaluation \\
Views & 17 rendered; 6 target views (front, front-right, right, back, left, front-left) used throughout \\
Scheduler / steps / seed & Euler discrete; 50 denoising steps; fixed seed 42 shared by all schedules within an experiment (identical initial latents) \\
Text prompt & None (reference-image-conditioned) \\
Appearance condition & Fixed reference rendering of the object (MVPainter's fixed reference-view input), identical across all compared schedules and methods \\
Geometry condition & 3-channel normal map $+$ 1-channel depth map $+$ 1-channel foreground mask \\
Metrics & FG-/Edge-SSIM, PSNR, FG-LPIPS (AlexNet) on foreground-masked regions; foreground Laplacian variance, RGB standard deviation, gradient magnitude on foreground pixels \\
Statistics & Two-sided paired tests; object-level percentile bootstrap 95\% CIs, 10{,}000 resamples; two-level (participant- and object-level) cluster bootstrap for the preference study \\
Code & Evaluation scripts (\texttt{geotex/}) and pipeline/configurations (\texttt{MVPainter/}) available with the checkpoints upon reasonable request (repository snapshot 194bce2) \\
\bottomrule
\end{tabular}
\end{table}

\end{document}
